\section{Discussion}
\label{sec:discussion}

The two factors that govern the benefit of self-supervision, the severity of missingness and the characteristics of the objective, admit a single interpretation: an objective helps to the extent that it supplies usable signal beyond what message passing already recovers. When observed features are abundant the encoder is not starved, so even a well-chosen objective adds little; when the target is class-irrelevant or trivially predictable, an objective adds little at any level of missingness. The largest gains therefore arise when both conditions align, that is, when a class-relevant, non-trivial objective is trained under severe missingness. This reading unifies the missingness-severity and objective-characteristic results under one mechanism rather than treating them as separate effects.

For practitioners, this suggests a concrete criterion for adding self-supervision under missing features: prefer objectives whose targets are class-relevant and not trivially predictable, expect the benefit to be largest under severe missingness and on dense continuous features, and expect little from self-supervision when features are abundant or sparse and binary. This stands in contrast to the common practice of importing an auxiliary objective from the graph self-supervision literature by analogy, without regard to whether its target carries downstream-relevant signal in the missing-feature regime. It also reconciles our findings with our earlier work on feature-based self-supervision for imputation: feature reconstruction remains useful, but because faithful reconstruction is neither necessary nor sufficient for an accurate decision boundary, a label-aware objective is the more reliable choice when the goal is classification rather than recovery.

It is also worth being precise about the source of NESS's overall strength. At moderate missingness, much of its performance derives from the encoder itself, namely the Feature-Propagation prefill together with a well-configured GraphSAGE backbone, while the self-supervised objectives act as a setting-dependent enhancement whose contribution grows as missingness becomes severe. Rather than the advantage resting on self-supervision alone, it is the combination that makes the model strong: the encoder and the auxiliary objectives are complementary, and together they add value across a range of settings and missing rates.

\paragraph{Limitations and future work.} Several limitations bound these conclusions and point to natural directions for future work. Our study considers missingness sampled uniformly at random, so whether the same criterion holds under structured or informative missingness remains open. We evaluate only on node classification, and extending the analysis to other downstream tasks, such as link prediction and regression, would test the generality of the class-relevance criterion. Our probe of target class-relevance is diagnostic after training rather than predictive before it; a measure of an objective's usefulness that can be computed in advance would turn the analysis into a design tool. Finally, the conclusion that neighborhood-centric self-supervision is inert on sparse binary attributes is drawn from small graphs, and a large sparse-attribute benchmark would strengthen it.
